\chapter{Silver's Theorem}

In order to prove Silver's Theorem, we will need a really powerful concept called stationary sets. Besides Silver's Theorem, we will also see stationary sets proves a lot of other important results.

\section{Clubs}

Before defining stationary sets, let's first analyze a similar concept.

Throughout this section, $\kappa$ will denote a regular cardinal.

\begin{definition}{Club}{club}
    A set $C\subseteq\kappa$ is called a \textcolor{laser_eyes!80!black}{closed unbounded set in $\kappa$} (or, more shortly, a \textcolor{laser_eyes!80!black}{club}), if
    \begin{itemize}
        \item $C$ is unbounded in $\kappa$;
        \item $C$ is closed, i.e., $\sup(C\cap\alpha)\in C$ for all $\alpha<\kappa$.
    \end{itemize}
\end{definition}

Note that $C$ being closed in $\kappa$ is equivalent to say that every sequence in $C$ of length $\alpha<\kappa$ has its limit in $C$. With this in mind, let's now see that clubs has the finite intersection property:

\begin{lemma}{}{club_fip}
    If $C_1$ and $C_2$ are clubs, then $C_1\cap C_2$ is also a club.
\end{lemma}

\begin{proof}
    To see it's closed, let $\alpha<\kappa$. Since $\sup(C_i\cap\alpha)\in C_i$, for $i=1,2$, then
    $$\sup((C_1\cap C_2)\cap\alpha)=\bigcup((C_1\cap\alpha)\cap(C_2\cap\alpha))=\rp{\bigcup(C_1\cap\alpha)}\cap\rp{\bigcup(C_2\cap\alpha)}\in C_1\cap C_2$$

    For unboundedness, let $\gamma<\kappa$, and consider the increasing sequence
    $$\gamma<\alpha_0<\beta_0<\alpha_1<\beta_1<\dots<\kappa$$
    where $\alpha_n\in C_1$ and $\beta_n\in C_2$, for $n\in\omega$. In fact, they exists, since $C_1$ and $C_2$ are unbounded, and their supremum $\delta$ is the same. Now, since $C_1$ and $C_2$ are also closed, we know $\delta\in C_1\cap C_2$.
\end{proof}

The previous lemma guarantees us the set of all clubs generates a filter, called \textcolor{laser_eyes!80!black}{the club filter} on $\kappa$
\begin{equation}\label{eq:club_filter}
    \mathrm{club}(\kappa):=\{X\subseteq\kappa:X\supseteq C\text{ for some club }C\}
\end{equation}
In fact, such filter is $\kappa$-complete, which follows from the next lemma

\begin{lemma}{}{club_k_complete}
    If $\{C_\alpha\}_{\alpha<\lambda}$ is a family of clubs, with $\lambda<\kappa$, then $\bigcap_{\alpha<\lambda}C_\alpha$ is also a club.
\end{lemma}

\begin{proof}
    The idea of the proof is similar to the one for two clubs.
    
    That it's closed is as easy as in the previous lemma. Given $\gamma<\kappa$, we have
    $$\bigcup\rp{\bigcap_{\alpha<\lambda}(C_\alpha\cap\gamma)}=\bigcap_{\alpha<\lambda}\rp{\bigcup(C_\alpha\cap\gamma)}\in\bigcap_{\alpha<\lambda}C_\alpha$$
    For unboundedness we can also try to follow a similar approach. In fact, as a warmup, consider $\kappa=\omega_1$ and $\lambda=\omega$. Then, given $\gamma<\omega_1$ we could consider the sequence
    $$\gamma<\alpha_0^0<\alpha_1^0<\alpha_1^1<\alpha_2^0<\alpha_2^1<\alpha_2^2<\alpha_3^0<\dots<\kappa$$
    where, for each $m\in\omega$, we have $\{\alpha_n^m\}_{n\in\omega}\subseteq C_m$.

    Now, in order to generalize this process, given $\gamma<\kappa$, let $\beta_0=\gamma$. Recursively, if $\beta_n$ is already defined, since each $C_\alpha$ is unbounded, we can find some $\delta_{\alpha,n}>\beta_n$ in $C_\alpha$. Let then $\beta_{n+1}=\sup_{\alpha<\lambda}\delta_{\alpha,n}$, which is a sequence of length $<\!\kappa$ with cardinals $<\!\kappa$. Since $\kappa$ is regular, $\beta_{n+1}<\kappa$, and then we can repeat the process, obtaining $(\beta_n)_{n\in\omega}$. Finally, let $\beta=\sup_{n\in\omega}\beta_n$, again, by regularity of $\kappa$ we have $\gamma<\beta<\kappa$ and since between each $\beta_n$, we have at least one element of each $C_\alpha$, it's clear that $\beta$ is also the supremum of each sequence $\delta_{\alpha,n}$ in $C_\alpha$, which is closed, so $\beta\in C_\alpha$ for each $\alpha<\lambda$, then
    $$\beta\in\bigcap_{\alpha<\lambda}C_\alpha$$
    as desired.
\end{proof}

\section{Stationary Sets and Fodor's Lemma}

\begin{definition}{Stationary Set}{stationary}
    A set $S\subseteq\kappa$ is \textcolor{laser_eyes!80!black}{stationary} if for every club $C$, we have $S\cap C\neq\emptyset$.
\end{definition}

By the previous lemmas, every club is stationary; and if $S$ is stationary, then any $T$ such that $S\subseteq T\subseteq\kappa$ is also stationary. Interestingly, nonstationary sets are exactly the ones belonging to the ideal dual to the club filter.

\begin{definition}{Regressive Function}{regressive}
    A function $f$ with domain $S\subseteq\kappa$ is \textcolor{laser_eyes!80!black}{regressive} if $f(\alpha)<\alpha$ for all $\alpha\neq0$.
\end{definition}

The next theorem\footnote{Which is usually called lemma} shows how uncountable regular cardinals differs substantially from $\aleph_0$. For example, it shows that there is no analogue for $\omega_1$ of the function $f:\omega\to\omega$ defined by $f(n)=n-1$ for $n>0$, and $f(0)=0$, which is regressive, and assume each value in $\omega$ only finitely many times.

You can to try to define it using Cantor Normal Form. For successor ordinals $\alpha+1$ it's easy, you can just let $f(\alpha+1)=\alpha$, but for limit ordinals, if you try, you will see you can't easily choose some smaller ordinal to be the image of $f$ without making it also the image of a lot of other ordinals.

You can't, for example, choose $f(\alpha)$, for $\alpha$ limit, to be ``the previous limit ordinal'', since, for example, $\alpha=\omega^2$ doesn't have one. You also can't just crop it by removing the last term in Cantor Normal Form, since then all ordinals of the form $\omega^\delta$ will be sent to $0$.

Anyway, it's a nice exercise to try.

\begin{theorem}{Fodor's Lemma}{fodor}
    A set $S\subseteq\kappa$ is stationary iff every regressive function $f:S\to\kappa$ is constant on an unbounded set. In fact, it will be constant in a stationary set.
\end{theorem}

Before proving it, we need an auxiliary lemma

\begin{lemma}{}{diagonal_intersection}
    If $\{C_\xi\}_{\xi<\kappa}$ is a family of clubs, then the set $C\subseteq\kappa$ defined by
    $$\alpha\in C\iff \alpha\in C_\xi,\;\forall \xi<\alpha$$
    called the \textcolor{laser_eyes!80!black}{Diagonal Intersection} of the $C_\xi$, is also a club.
\end{lemma}

\begin{proof}
    To see it's closed, let $\lambda<\kappa$ and $(\alpha_\beta)_{\beta<\lambda}$ be an increasing sequence in $C$. Now, if $\xi<\alpha=\sup\{\alpha_\beta\}_{\beta<\lambda}$, by definition $\xi<\alpha_\delta$ for some $\delta<\lambda$. But then $\xi<\alpha_\beta$, for all $\beta\geq\delta$. Since each $\alpha_\beta\in C$, then $\alpha_\beta\in C_\xi$, for all $\beta\geq\delta$, but since $C_\xi$ is closed, $\alpha=\sup\{\alpha_\beta:\beta\geq\delta\}\in C_\xi$ too. So we proved that $\alpha\in C_\xi$ for each $\xi<\alpha$, i.e., $\alpha\in C$.

    For unboundedness, fix $\gamma<\kappa$, let's construct an increasing sequence $(\alpha_n)_{n\in\omega}$ as recursively as follows: let $\alpha_0=\gamma$ and, if $\alpha_n<\kappa$ is already defined, we know by \cref{lem:club_k_complete} that $\bigcap_{\xi<\alpha_0}C_\xi$ is a club, then we can let $\alpha_{n+1}$ be the least element in it greater than $\alpha_n$.

    Now, let $\alpha=\sup_{n\in\omega}\alpha_n$, we claim $\alpha\in C$. In fact, if $\xi<\alpha$, by definition $\xi<\alpha_k$ for some $k\in\omega$, but then, by construction, $\alpha_{n+1}\in\bigcap_{\beta<\alpha_n}C_\beta$, in particular, $\alpha_{n+1}\in C_\xi$, for all $n\geq k$.
\end{proof}

We can now prove Fodor's Lemma.

\begin{proof}
    $(\Leftarrow)$ Let $S\subseteq\kappa$ be nonstationary, it's sufficient to exhibit an example of a regressive $f:S\to\kappa$ that takes each value $<\!\kappa$-many times. Since $S$ is nonstationary, there is some club $C$ disjoint from $S$. Let then, for each $\alpha\in\kappa$, $f$ be define as follows:
    $$f(\alpha)=\sup(C\cap\alpha)$$
    Since $f(\alpha)\in C$, $\alpha\in S$, and $S\cap C=\emptyset$, then $f(\alpha)<\alpha$. However, given any $\gamma<\kappa$, if $\delta$ is the least element of $C$ greater than $\gamma$, then for any $\alpha>\delta$ we have
    $$f(\alpha)=\sup(C\cap\alpha)\geq \delta>\gamma$$
    and then it can't be constant in $\gamma$ on an unbounded set.

    $(\Rightarrow)$ For each $\xi<\kappa$, let $S_\xi=f^{-1}(\xi)$, we will show $S_\xi$ is stationary for some $\xi$.

    Otherwise, for each $\xi<\kappa$ there would be a club $C_\xi$ disjoint from $S_\xi$. Let $C$ be the diagonal intersection of the $C_\xi$. Then $\alpha\in C$ iff $\alpha\in C_\xi$ for all $\xi<\alpha$, i.e., $f(\alpha)\neq\xi$ for all $\xi<\alpha$, this basically says that $f(\alpha)\geq\alpha$. But, by the previous Lemma, $C$ is a club, hence it intersects $S$; however, $f$ is regressive in $S$, a contradiction.
\end{proof}

\subsection{Applications}

Let's now see how powerful Fodor's Lemma is.

First note that, until now, we saw every club is stationary, and then that every superset of a club is also stationary. However, we didn't see any example of a stationary set that doesn't contain a club. The next corollary will show us we can in fact find one\footnote{It's worth noting we can only prove it under the Axiom of Choice, it's consistent with ZF that every stationary set contains a club.}

\begin{corollary}{}{complement_stationary}
    There are arbitrarily large stationary sets whose complement is also stationary.
\end{corollary}

\begin{proof}
    Let $C$ be the set of all limit ordinals below $\kappa$. For each $\alpha\in C$, there exists an increasing sequence $(x_{\alpha,n})\to \alpha$. For each $n\in\omega$, let $f_n:C\to\kappa$ be given by $f_n(\alpha)=x_{\alpha,n}$.

    Since $f_n(\alpha)<\alpha$ for every $n\in\omega$, by a Scholium of Fodor's Lemma to each $f_n$, we know there exists $\gamma_n$ such that the set $S_n=f^{-1}_n(\gamma_n)$ is stationary, for each $n\in\omega$.

    We claim at least one of the $S_n$'s has a stationary complement. If not, each one contains a club, hence $\bigcap_{n\in\omega}S_n$ contains their intersection, which is also a club. By unboundedness it has some ordinal $\alpha>\sup_{n\in\omega}\gamma_n$, but then $x_{\alpha,n}=f_n(\alpha)=\gamma_n$ doesn't converge to $\alpha$, a contradiction.
\end{proof}

As another application, we will prove a really famous theorem\footnote{Although it's called lemma, which we will in fact state as a corollary xD.} in combinatorial set theory, known as the $\Delta$-System Lemma. In particular, it's really important when proving some Forcing Notions are in fact ccc. Let's first define a $\Delta$-system.

\begin{definition}{}{delta_system}
    A family of sets $\mc{A}$ forms a \textcolor{laser_eyes!80!black}{$\Delta$-system with root $R$} if, and only if, $X\cap Y=R$, whenever $X,Y\in\mc{A}$, with $X\neq Y$.
\end{definition}

We also say $\mc{B}$ forms a $\Delta$-system, if there is a set $R$ such that $\mc{B}$ is a $\Delta$-system with root $R$.

\begin{corollary}{$\Delta$-System Lemma}{delta_system}
    Let $\mc{A}$ be a family of finite sets with $\lvert\mc{A}\rvert=\kappa$. Then there is a $\mc{B}\in[\mc{A}]^\kappa$ such that $\mc{B}$ forms a $\Delta$-system.
\end{corollary}

\begin{proof}
    We may assume WLOG that $\mc{A}\subseteq[\kappa]^{<\omega}$. Since every set in $\mc{A}$ is finite, and we have $\kappa$-many of those, by Pigeonhole we can find $\kappa$ of them $\{A_\alpha:\alpha<\kappa\}$ of the same size $n\in\omega$\todo{todo}...
\end{proof}

\section{Silver's Theorem}

Let's now state the main theorem of this section. Roughly, it says that if GCH holds for all cardinals smaller than some singular with uncountable cofinality, then it also holds for it.

\begin{theorem}{Silver's Theorem}{silver}
    Let $\aleph_\lambda$ be a singular cardinal, with $\cf(\lambda)>\omega$. If $2^{\aleph_\alpha}=\aleph_{\alpha+1}$, for every $\alpha<\lambda$, then
    $$2^{\aleph_\lambda}=\aleph_{\lambda+1}$$
\end{theorem}

\subsection{Original Proof Sketch}

The proof of Silver's theorem is quite intricate, and perhaps many of the steps and lemmas do not seem natural. In fact, it is worth mentioning that the original proof is much more complicated, involving, for example, the use of forcing.

The original theorem appeared in 1974 at \todo{add reference}..., in an article called ``On the singular cardinals problem'', and stated that GCH couldn't fail for the first time in a singular cardinal of uncountable cofinality. This result was inspired by the work of Magidor, who at the time was interested in proving that certain types of ultrafilters, called regular ultrafilters, did not exist in $\omega_1$. Instead, he proved that if they did exist, Silver's theorem holds for $\lambda=\omega_1$!\footnote{Not a factorial joke.}

So, roughly, based on Magidor's results, instead of trying to construct directly such object, Silver extends the universe through forcing! He adds to the ground model $M$ a bijection between $2^{\cf(\kappa)}$ and $\aleph_0$ to make it countable, which simplify the combinatorics a lot. He then construct an ultraproduct of the ground model in which $\mc{P}(\kappa)$ is small, allowing him to ``reflect'' such result and conclude $\mc{P}(\kappa)$ is also small in the ground model $M$. A complete proof can be found in \todo{add reference}...

\subsection{Proof}
In order to prove it, we will first prove a series of lemmas needed to show Silver's result.

\begin{definition}{Strong Limit Cardinal}{strong_limit}
    A cardinal $\kappa$ is called a \textcolor{laser_eyes!80!black}{strong limit cardinal}\footnote{Yeah, this is the same strong limit we cited in the end of \cref{sec:weakly_inaccessible}} if, for any $\lambda<\kappa$, we have $2^\lambda<\kappa$.
\end{definition}

It's easy to see that, under GCH, the concept of a limit cardinal matches with the one of a strong limit cardinal. Why we should care? Look again to Silver's Theorem hypothesis, we have some kind of partial GCH for all cardinal smaller than $\aleph_\lambda$. In fact, this imply that $\aleph_\alpha$ is a strong limit cardinal for every $\alpha<\lambda$. In particular, the next lemma holds.

\begin{lemma}{}{strong_limit_exponentiation}
    If $\aleph_\alpha$ is a strong limit cardinal, and if $\kappa,\lambda<\aleph_\alpha$ are infinite cardinals, then $\kappa^\lambda<\aleph_\alpha$.
\end{lemma}

\begin{proof}
    $$\kappa^\lambda\leq(\kappa\cdot\lambda)^{\kappa\cdot\lambda}=2^{\kappa\cdot\lambda}<\aleph_\alpha$$
    where in the inequality we used \cref{lem:card_exp_a_leq_b}.
\end{proof}

\textcolor{laser_eyes!80!black}{Throughout this section, we will fix a $\lambda$ as in Silver's Theorem, and let $\kappa=\cf(\lambda)<\lambda$}. Then, by the previous Lemma, we can conclude the following, which we will use extensively from now on

\begin{equation}\label{eq:silver_aux_eq}
    2^\kappa<\aleph_\lambda\text{ and }\aleph_\alpha^\kappa<\aleph_\lambda\text{, for every }\alpha<\lambda.
\end{equation}

\begin{definition}{}{almost_disjoint}
    Two functions $f,g$ on $\kappa$ are \textcolor{laser_eyes!80!black}{almost disjoint} if there is some $\alpha<\kappa$ such that $f(\beta)\neq g(\beta)$, for all $\beta\geq\alpha$.
\end{definition}

We can now state our second lemma.

\begin{lemma}{}{fodor_aux1}
    Let $\{A_\alpha:\alpha<\kappa\}$ be a family such that $T=\{\alpha<\kappa:\lvert A_\alpha\rvert\leq\aleph_\alpha\}$ is stationary. If $F$ is a collection of almost disjoint choice functions for the family, then $\lvert F\rvert\leq\aleph_\lambda$.
\end{lemma}

\begin{proof}
    We can assume WLOG that $A_\alpha\subseteq\omega_\alpha$. Let $S_0$ be the set of all nonzero limit ordinals in $T$. Given $f\in F$ and $\alpha\in S_0$, let $g(\alpha)$ be the least $\beta$ such that $f(\alpha)<\omega_\beta$.

    Since $f(\alpha)\in A_\alpha\subseteq\omega_\alpha$, we have $f(\alpha)<\omega_\alpha$ and, since $\alpha$ is a limit ordinal, then $f(\alpha)<\omega_\beta$, for some $\beta<\alpha$. This implies that $g(\alpha)<\alpha$, i.e., it's a regressive function. Since $S_0$ is stationary, by Fodor's Lemma there is some stationary $S\subseteq S_0$ such that $g\restriction_S$ is constant, let's say in $\beta$. Then $f\restriction_S$ is a function from $S\subseteq\kappa$ to $\omega_\beta<\omega_\kappa$. Note that both $S$ and $\beta$ depends on $f$.

    Now, consider the function $\varphi$ that takes some $f\in F$ and return exactly that $f\restriction_S:S\to\omega_\beta$. Then
    $$\varphi:F\to\bigcup_{S\subseteq\kappa}\rp{\bigcup_{\beta<\kappa}\omega_\beta^S}$$
    We claim $\varphi$ is injective. In fact, if $f\neq g$, even if $\varphi(f)$ and $\varphi(g)$ have the same domain $S$, since they're almost disjoint, there is some $\alpha<\kappa$ such that $f(\beta)\neq g(\beta)$ for every $\beta\geq\alpha$. But since $S$ is stationary, there is some $\gamma>\alpha$ in it, so $\varphi(f)(\gamma)\neq\varphi(g)(\gamma)$. Then, since $\varphi$ is injective,
    $$\lvert F\rvert\leq\sum_{S\subseteq\kappa}\sum_{\beta<\kappa}\aleph_\beta^{\lvert S\rvert}\leq 2^\kappa\sum_{\beta<\kappa}\aleph_\beta^\kappa\leq\aleph_\lambda$$
    where in the last inequality we used \cref{eq:silver_aux_eq}.
\end{proof}

We can now state our main lemma to prove Silver's Theorem.

\begin{lemma}{}{fodor_aux2}
    Let $\{A_\alpha:\alpha<\kappa\}$ be a family such that $T=\{\alpha<\kappa:\lvert A_\alpha\rvert\leq\aleph_{\alpha+1}\}$ is stationary. If $F$ is a collection of almost disjoint choice functions for the family, then $\lvert F\rvert\leq\aleph_{\lambda+1}$.
\end{lemma}

\begin{proof}
    Let $U$ be an ultrafilter extending the club filter (as defined in (\ref{eq:club_filter})). As in the previous lemma, let's assume $A_\alpha\subseteq\omega_{\alpha+1}$, and let's define the relation $<$ on $F$ induced by $U$ as follows:
    $$f<g\iff \{\alpha<\kappa:f(\alpha)<g(\alpha)\}\in U$$
    We claim $(F,<)$ is a linear order. In fact, if $f<g$ and $g<h$, then
    $$\{\alpha<\kappa:f(\alpha)<h(\alpha)\}\supseteq\{\alpha<\kappa:f(\alpha)<g(\alpha)\}\cap\{\alpha<\kappa:g(\alpha)<h(\alpha)\}$$
    Since $U$ is a filter, the RHS is in it, and since it's complete, and we can't have the LHS outside of it, then it's also in $U$, so $f<h$ and $<$ is transitive.

    Now, if $f\neq g$, since they're almost disjoint we have $\{\alpha<\kappa:f(\alpha)=g(\alpha)\}$ bounded, so it's not in $U$, meaning that either $\{\alpha<\kappa:f(\alpha)<g(\alpha)\}$ or $\{\alpha<\kappa:g(\alpha)<f(\alpha)\}$ is in $U$, i.e., $<$ satisfies trichotomy. This concludes it's a linear ordering on $F$.

    Now, note that every $S\in U$ is stationary, since otherwise $S\cap C=\emptyset$ would be in $U$ for some club $C$. Hence, if $f,g\in F$ and $g<f$, then $g\in F_f$, where
    $$F_f=\{g\in F:\text{for some stationary }T,\; g(\alpha)<f(\alpha)\text{ for all }\alpha\in T\}$$
    Note that we have
    $$F_f\subseteq\bigcup_{T\subseteq\kappa}\{g\in F:g(\alpha)<f(\alpha)\text{ for all }\alpha\in T\}$$
    let's estimate the RHS. In fact, each $\{g\in F:g(\alpha)<f(\alpha)\text{ for all }\alpha\in T\}\subseteq\prod_{\alpha<\kappa}B_\alpha$ is such that $g(\alpha)<f(\alpha)<\omega_{\alpha+1}$ for $\alpha\in T$, then we can let $B_\alpha=f(\alpha)$ and then $\lvert B_\alpha\rvert\leq\aleph_\alpha$, i.e., we're under the hypothesis of the previous lemma. So we can conclude that each term in the union has size at most $\aleph_\lambda$, therefore
    $$\lvert F_f\rvert\leq\sum_{T\subseteq\kappa}\aleph_\lambda=2^\kappa\cdot \aleph_\lambda=\aleph_\lambda$$
    where, again, we used \cref{eq:silver_aux_eq}. We then have $(F,<)$ a linear order whose initial segments has all cardinality at most $\aleph_\lambda$, which implies\footnote{In general, under AC, if $\lvert\mathrm{pred}(a)\rvert<\aleph_\gamma$, for every $a$ in a linear ordered set $A$, then $\lvert A\rvert\leq\aleph_\gamma$.} $\lvert F\rvert\leq\aleph_{\lambda+1}$.
\end{proof}

And now we can finally prove Silver's Theorem.

\begin{proof}
    Let $A_\alpha=\mc{P}(\omega_\alpha)$, for $\alpha<\kappa$. By hypothesis $\lvert A_\alpha\rvert=2^{\aleph_\alpha}=\aleph_{\alpha+1}$. Now, for each $X\subseteq\omega_\lambda$, let $f_X$ be a choice function defined in the family $\{A_\alpha:\alpha<\kappa\}$ as follows
    $$f_X(\alpha)=X\cap\omega_\alpha$$
    In fact, if $X\neq Y$, we can get the witnesses at some stage $\omega_\alpha$, and then $f_X(\beta)\neq f_Y(\beta)$ for all $\beta\geq\alpha$, then $F=\{f_X:X\in\mc{P}(\omega_\lambda)\}$ is a family of almost disjoint functions. By the previous lemma
    $$\aleph_\lambda<2^{\aleph_\lambda}=\lvert F\rvert\leq\aleph_{\lambda+1}\implies 2^{\aleph_\lambda}=\aleph_{\lambda+1}$$
\end{proof}

\section{Singular Cardinal Hypothesis}

